\chapter{Conclusion}
\label{ch:conclusion}

This thesis investigated whether recent vision and multimodal foundation models can support the quality assessment and selective refinement of existing building-footprint datasets in regions where authoritative reference data are often limited or unavailable. Instead of generating a new building dataset from scratch, the work focused on an update-oriented workflow: existing Google Open Buildings polygons were visually assessed against high-resolution OpenAerialMap imagery and, where necessary, refined using MLLM-guided prompting, SAM 3-based segmentation, and geometric post-processing.

The results show that this type of workflow can meaningfully improve the visual agreement between existing building polygons and the evaluated imagery. Across the five selected study areas, the workflow was applied to 2,842 Google Open Buildings polygons. Of these, 1,559 were routed as visible partial- or full-house cases, 1,518 produced a SAM-derived polygon, and 1,352 resulted in a post-processed output. The final manual evaluation included 1,256 buildings. Within this subset, 82.8\% of the evaluated polygons improved after post-processing, 13.1\% remained largely unchanged, and 4.1\% degraded. The share of polygons rated as \textit{good} or \textit{perfect} increased substantially after refinement, indicating that the proposed workflow was able to correct many visibly problematic building geometries.

With respect to the first research question, the results suggest that an MLLM can support visual quality assessment of existing vector-based building footprints, but mainly at a coarse decision-making level. The Qwen3 8B model was useful for identifying whether a building was visible, deciding whether a case should be processed as a full-house or partial-house example, and generating approximate positive and negative prompt points for SAM-based segmentation. The point-prompt evaluation showed that the spatial prompt generation worked reasonably well in successful full-house cases, with more than 90\% of positive and negative points being spatially consistent with the final building polygons. However, the MLLM was not reliable enough for fine-grained semantic error description. It struggled to consistently distinguish between shifted polygons, missing parts, extra parts, oversimplified outlines, and general shape mismatches. Therefore, the MLLM component should currently be understood as a useful routing and prompting aid, but not as a complete replacement for deterministic quality metrics or expert interpretation.

The second research question concerned the effectiveness of MLLM-generated prompts, SAM-based segmentation, and post-processing for improving selected building-footprint geometries. The strongest improvement was observed for positional displacement. The share of polygons classified as \texttt{SHIFTED} decreased from 77.3\% in the original Google Open Buildings polygons to 4.3\% after post-processing. \texttt{SHAPE\_MISMATCH} and \texttt{OVERSIMPLIFIED} errors were also reduced substantially. This demonstrates that the workflow was particularly effective for cases where an existing polygon already indicated the approximate presence of a building but did not align well with the visible roof structure. At the same time, the results also show that refinement introduced new challenges. The categories \texttt{MISSING\_PARTS} and \texttt{EXTRA\_PARTS} increased after processing, and new geometric errors were introduced in 16.5\% of the manually evaluated buildings. These errors were often related to ambiguous segmentation masks, vegetation, occlusion, complex roof structures, or the trade-off between preserving visual detail and producing clean GIS-ready vector geometries.

The third research question addressed transferability across different geographic regions, building morphologies, and imagery conditions. The workflow improved the majority of evaluated polygons in all five study areas, but the degree of improvement varied substantially. Niger and Mozambique achieved the strongest overall results, while Liberia produced strong but partly development-influenced results. Mexico and Nepal were more difficult cases, with lower final quality, larger area changes, and stronger increases in polygon complexity. These differences show that transferability is possible, but not uniform. The workflow performed best where buildings were visually separable and morphologically compact. It performed less reliably in areas with incremental housing, heterogeneous roof materials, dense structural complexity, vegetation, or ambiguous object boundaries. The supplementary Switzerland test further showed that the workflow is not automatically beneficial for already high-quality building datasets. In that case, the regularized outputs performed slightly worse than the existing Microsoft source polygons when compared against authoritative TLM geometries. This supports the interpretation that the method is most useful as a selective correction workflow for visibly problematic polygons, rather than as a universal improvement tool.

A central limitation of the thesis is that the evaluation was image-relative. The OpenAerialMap TIFF imagery used in this work was not necessarily identical to the imagery used to generate the original Google Open Buildings dataset. Therefore, the corrected polygons should be interpreted as geometrically improved with respect to the evaluated imagery, not necessarily as absolutely more accurate representations of the real-world building footprint at the time of analysis. This distinction is important because temporal differences, georeferencing uncertainty, image registration effects, and orthorectification differences may all contribute to apparent polygon-image mismatches.

Overall, the thesis demonstrates that multimodal and visual foundation models can contribute to local building-footprint quality assessment and refinement, but that they do not yet provide a fully automatic and reliable geospatial QA solution. The main contribution lies in showing how an existing global building dataset can be combined with MLLM-based visual reasoning, prompt generation, SAM-based segmentation, and vector post-processing in a modular workflow. The results indicate that such systems are especially promising when used in a human-in-the-loop setting: they can identify problematic polygons, generate candidate corrections, and support local mapping or database maintenance, while final decisions remain with human experts.

Future development should focus on more robust building regularization, better handling of occlusion and vegetation, systematic fine-tuning or evaluation of stronger spatial MLLMs, and closer integration into practical mapping workflows. If more capable geospatial foundation models become available, the MLLM component may improve substantially, especially for building detection, prompt placement, and semantic error interpretation. However, even with stronger models, reliable geospatial data maintenance will likely require a combination of visual AI, deterministic geometry checks, uncertainty communication, and human validation. In this sense, the workflow developed in this thesis should be understood not as a replacement for expert mapping, but as a step toward collaborative, AI-supported geospatial quality assessment for large-scale building-footprint datasets.

